\section{实验一：开发环境配置}

\subsection{实验目的}

本实验的目标是完成机械臂综合实验的软件环境整理，为后续运动控制、视觉感知和抓取程序提供统一的仿真与代码运行基础。环境配置包括 Python 依赖、机械臂串口通信方式说明、仿真模型文件、实验代码目录和安全运行约束；本组未连接真实机械臂验证硬件通信。

\subsection{软件环境}

本实验仓库根目录提供了 \texttt{environment.yml}，用于创建名为 \texttt{robot} 的 conda 环境。环境中 Python 版本为 3.10.20，主要依赖包括：

\begin{table}[H]
  \centering
  \caption{实验环境主要软件依赖}
  \small
  \begin{tabularx}{\textwidth}{L{2.3cm}Y Y}
    \toprule
    类别 & 软件包 & 用途 \\
    \midrule
    机器人建模 & \texttt{roboticstoolbox-python}, \texttt{spatialmath-python} & DH 建模、正逆运动学、轨迹规划 \\
    数值计算 & \texttt{numpy}, \texttt{scipy}, \texttt{modern-robotics} & 矩阵计算、优化求解 \\
    视觉处理 & \texttt{opencv-python}, \texttt{opencv-contrib-python} & 相机标定、颜色分割、图像显示 \\
    仿真与显示 & \texttt{matplotlib}, \texttt{pybullet}, \texttt{trimesh} & 轨迹可视化、模型检查 \\
    深度学习 & \texttt{torch}, \texttt{torchvision} & YOLO 训练和推理基础 \\
    通信控制 & \texttt{pyserial} & 串口转 CAN 通信 \\
    \bottomrule
  \end{tabularx}
\end{table}

首次配置时可在仓库根目录执行：

\begin{lstlisting}[language=bash,caption={conda 环境创建命令}]
conda env create -f environment.yml
conda activate robot
\end{lstlisting}

后续运行实验脚本时，应优先在 \texttt{daran/} 或其子目录下执行，以保证相对路径中的模型文件、字体文件和标定结果可以正确加载。

\subsection{环境配置原则}

机器人综合实验的环境配置与普通 Python 项目不同，原因在于本实验同时依赖数值计算库、机器人运动学库、相机视觉库、图形界面库以及真实串口设备。若某一类依赖版本不一致，问题往往不会立刻表现为导入失败，而是表现为仿真窗口无法打开、相机标定结果无法保存、机械臂模型姿态异常或真实关节无法通信。因此本实验采用“一套环境、一套模型、一套路径约定”的配置原则。

第一，所有 Python 代码统一运行在 \texttt{robot} 环境中。该环境同时包含 \texttt{roboticstoolbox-python}、\texttt{spatialmath-python}、\texttt{opencv-contrib-python}、\texttt{pyserial} 和 \texttt{torch} 等包，能够覆盖从 DH 建模、轨迹规划到视觉识别的主要需求。第二，所有实验脚本尽量从 \texttt{daran/} 内部运行，因为机械臂模型、字体、标定结果和图像数据大量使用相对路径。第三，涉及真实机械臂的程序必须显式区分仿真模式和实机模式，避免在检查代码或生成图像时误触发硬件动作。

从依赖来源看，\texttt{environment.yml} 使用清华镜像和 conda-forge 作为主要 channel，并将部分机器人库和视觉库通过 pip 安装。这样做可以减少课程实验中常见的二进制包冲突问题。尤其是 \texttt{roboticstoolbox-python} 依赖的 \texttt{spatialgeometry}、\texttt{swift-sim}、\texttt{rtb-data} 等包需要保持兼容，否则 DH 模型可视化和轨迹动画可能出现运行时错误。

\subsection{硬件连接与安全约束}

机械臂通过串口转 CAN 模块与上位机通信，默认波特率为 115200。代码中常见串口名称包括 Windows 下的 \texttt{COM4}/\texttt{COM7}，Linux 或 Jetson 下的 \path{/dev/ttyUSB0}，以及部分脚本中使用的 \path{/dev/serial/by-id/...}。由于 \texttt{arm\_robot} 对象在实例化后会打开串口并控制真实硬件，实验中必须区分仿真脚本和实机脚本。

本仓库采用“先仿真、后实机”的安全模式。以任务 1 和任务 2 为例，脚本中使用 \texttt{MOVE\_REAL\_ROBOT} 或人工确认按钮控制是否下发真实动作；在未来打开实机开关前，需要先检查关节限位、桌面高度、目标点可达性和仿真动画。本报告仅分析该安全设计并运行仿真，不进入实机状态。

从实验安全角度看，环境配置并不是简单的软件安装步骤，而是后续所有实验的前置保护条件。若 Python 环境、串口设备名、机械臂模型参数或标定文件路径不一致，同一段运动程序可能出现仿真正常而实机异常的情况。因此本实验首先确认依赖版本、目录结构和执行入口，并明确区分仿真验证与真实指令下发；真实硬件链路留待后续单独验证。

代码中的真实机械臂控制链路可以分为四层：最底层为串口转 CAN 通信，中间层为单关节角度与属性读写，上层为六轴机械臂的正逆运动学和夹爪控制，最外层为实验脚本中的任务流程。本次仅确认相关模块和接口的代码组织，没有验证串口设备、关节反馈或真实动作。后续若开展实机实验，应先读取关节角、显示仿真姿态并发送低速小幅目标确认通信状态，再考虑执行完整任务。

\subsection{仓库结构与实验代码组织}

实验代码主要位于 \texttt{daran/} 目录下。核心机械臂控制代码采用逐层继承结构：底层 \texttt{DrEmpower\_can.py} 负责串口/CAN 协议，\texttt{arm\_six\_axis.py} 实现六轴运动学，\texttt{gripper.py} 添加夹爪控制，\texttt{arm\_robot.py} 提供面向实验脚本的机械臂控制接口。

上述继承关系决定了后续实验的代码阅读顺序。若只关注实验脚本，可以从 \texttt{arm\_robot} 的 \texttt{set\_arm\_joints}、\texttt{read\_joints} 和夹爪接口入手；若需要解释运动学，则应回到 \texttt{arm\_six\_axis.py} 或 DH notebook；若需要分析真实执行异常，则要查看 \texttt{DrEmpower\_can.py} 中的串口通信和关节属性轮询。报告正文主要讨论实验层面的流程，但在解释误差来源时，需要同时考虑底层通信、运动学模型和执行机构三方面因素。

与本报告六个实验对应的主要文件如表 \ref{tab:repo-map} 所示。

\begin{table}[H]
  \centering
  \caption{实验内容与仓库文件对应关系}
  \label{tab:repo-map}
  \small
  \begin{tabularx}{\textwidth}{L{2.2cm}Y Y}
    \toprule
    实验 & 主要目录或文件 & 内容 \\
    \midrule
    环境配置 & \path{environment.yml}, \path{CLAUDE.md} & 环境、依赖、运行约束 \\
    直线运动 & \path{daran/task1/} & DH 建模、IK、直线轨迹仿真 \\
    抓取搬运 & \path{daran/task2/pick_and_place.py} & 示教点回放和夹爪控制 \\
    相机标定 & \path{daran/task3/camera_calibration.py} & 棋盘格内参标定 \\
    颜色视觉抓取 & \path{daran/task3/visual_grasp.py}, \path{daran/color_pick_place/} & 颜色检测、坐标转换、抓取规划 \\
    YOLO 视觉抓取 & \path{daran/color_pick_place/yolo/} & 数据集、训练、推理与接口复用 \\
    \bottomrule
  \end{tabularx}
\end{table}

\subsection{配置验证步骤}

环境配置完成后，本次从软件依赖、模型加载和运行安全三个层面进行检查。首先进入 \texttt{robot} 环境，检查 \texttt{numpy}、\texttt{roboticstoolbox}、\texttt{cv2}、\texttt{torch} 等核心包；其次运行只涉及模型和图形显示的脚本，检查 DH 模型、关节限位和仿真输出；最后阅读涉及串口控制的脚本，确认真实执行需要人工切换或按钮确认，但不连接硬件。

本实验采用的流程是“先静态检查，再仿真运行，硬件连接留待后续”。静态检查用于发现缺包、路径错误和配置文件缺失；仿真运行用于检查机械臂模型与轨迹生成逻辑。本组没有进行串口可用性、通信参数或真实关节动作验证，因此相关内容只能作为后续实机测试方案。

配置验证还包括对路径和输出文件的检查。例如任务 1 的仿真图片位于 \path{daran/task1/}，任务 2 的示教点位于 \path{daran/task2/ref_pose.json}，任务 3 的相机标定结果位于 \path{daran/task3/calibration_output/}。这些文件不仅是程序运行产物，也是报告撰写中的实验依据。如果路径不一致，后续章节中引用的图像、表格和误差指标都会失去来源。

\begin{table}[H]
  \centering
  \caption{环境配置后的验证项目}
  \label{tab:env-checklist}
  \small
  \begin{tabularx}{\textwidth}{L{2.6cm}Y Y}
    \toprule
    验证项目 & 检查内容 & 通过标准 \\
    \midrule
    Python 环境 & \texttt{robot} 环境、Python 3.10、核心包导入 & 依赖无缺失，导入不报错 \\
    DH 模型 & \texttt{RevoluteDH} 参数、关节限位、零位和 park 姿态 & 模型姿态与预期一致，坐标轴显示正常 \\
    图形输出 & matplotlib/PyPlot 后端、中文字体、GIF 或 PNG 保存 & 仿真图像能生成并保存到实验目录 \\
    相机与图像 & OpenCV 读取图像、棋盘格检测、标定文件保存 & 能检测角点并生成 JSON/NPZ 结果 \\
    串口通信 & 设备端口、波特率、关节读取接口 & 本次未连接硬件，后续需单独验证 \\
    \bottomrule
  \end{tabularx}
\end{table}

\subsection{环境复现与常见问题}

为了使报告中的实验能够被复现，环境记录不能只写“安装 Python 和 OpenCV”。本实验中较关键的版本包括 Python 3.10.20、\texttt{numpy} 1.26.4、\texttt{opencv-python} 4.11.0.86、\texttt{roboticstoolbox-python} 1.1.1、\texttt{spatialmath-python} 1.1.16、\texttt{pyserial} 3.5 和 \texttt{torch} 2.11.0。机器人建模、相机标定和 YOLO 推理分别依赖这些包的不同部分，因此版本漂移可能在不同实验中表现出不同问题。

实际配置中最常见的问题有三类。第一类是图形后端问题，例如 PyQt 与系统 Qt 库版本冲突，表现为 GUI 无法启动或出现 \texttt{xcb} 相关错误。任务 1 的 GUI 程序中加入了清理系统 \texttt{LD\_LIBRARY\_PATH} 后重启进程的逻辑，就是为了解决这类问题。第二类是路径问题，例如从仓库根目录运行脚本时找不到 \texttt{arm\_robot.py}、字体或标定文件。第三类是真实设备问题，例如串口名称变化、权限不足或 CAN 设备未连接。前两类问题可以通过静态检查和仿真验证发现，第三类必须在实机连接阶段谨慎确认。

在课程实验报告中写清这些问题有两个作用。一方面，它能解释为什么同一套代码在不同机器上可能出现不同表现；另一方面，它能说明实验设计中为什么反复强调 dry-run、仿真预览和人工确认。机器人实验的可靠性不只来自算法正确，也来自运行环境、执行入口和安全流程的一致。

\subsection{实验结果与小结}

通过环境配置实验，整理了 conda 环境、核心依赖、机械臂控制接口和实验目录结构。后续程序均基于同一个 \texttt{robot} 环境和同一套机械臂模型参数进行代码分析与仿真。安全方面，本报告只讨论不触发真实硬件动作的仿真和程序逻辑；涉及真实机械臂的脚本仅作为实现材料，未验证其硬件运行效果。

本实验也明确了后续章节的共同约束：运动控制程序应先在 DH 模型中检查关节限位和桌面高度；抓取搬运程序应先检查参考点的安全性；视觉程序应先确认相机内参、外参和坐标转换文件存在。也就是说，环境配置不是孤立的第一步，而是整套仿真与代码验证流程的入口；实机效果仍需在硬件条件具备后另行测试。
